% Бүлэг 4

\chapter{Эдийн засгийн үр ашгийн тооцоо}
\label{Chapter4}

\section{Төслийн эдийн засгийн үр ашгийн шинжилгээ}
Мэдээллийн системийн төслийг зөвхөн техникийн хэрэгжилтээр бус, эдийн засгийн үр ашгийн үзүүлэлтээр үнэлэх шаардлагатай. Энэхүү систем нь байгууллагын дотоод хэрэгсэл тул түүний өгөөж нь борлуулалтын орлого бус, ажилтнуудын хэмнэсэн ажлын цагийн үнэлгээ юм. Иймд үр ашгийг 25 ажилтантай (үүнээс 5 менежер) жишээ байгууллагын хэмжээнд, мөнгөний цаг хугацааны үнэ цэнийг тооцсон аргуудаар (NPV, нөхөх хугацаа) болон ашигт ажиллагаа, зардлын түвшин, бүтээмжийн үзүүлэлтүүдээр тооцов.

\subsection{Тооцоололд ашигласан томьёонууд}
Мөнгөний өнөөгийн үнэ цэнэ (PV):
\begin{equation}
PV = \frac{FV}{(1+r)^n}
\end{equation}
Хөрөнгө оруулалтын өнөөгийн цэвэр үнэ цэнэ (NPV):
\begin{equation}
NPV = \sum_{i=1}^{n} \frac{FV_i}{(1+r)^i} - Inv
\end{equation}
энд $FV_i$ -- $i$-р жилийн цэвэр мөнгөн урсгал, $r$ -- хорогдлын хувь, $n$ -- хугацаа, $Inv$ -- анхны хөрөнгө оруулалт.

Хөрөнгө оруулалтыг нөхөх хугацаа:
\begin{equation}
T_{inv} = t + \frac{Inv - \sum_{i=1}^{t} PV_i}{PV_{t+1}}
\end{equation}
Хөрөнгө оруулалтын үр ашгийн дундаж норм:
\begin{equation}
PfN = \frac{\left(\sum_{i=1}^{n} PV_i\right)/n}{Inv}
\end{equation}
Ашигт ажиллагаа, зардлын түвшин, нийт бүтээмж:
\begin{equation}
P_a = \frac{Profit}{Income}\times 100\%, \quad
P_c = \frac{Profit}{Inv}, \quad
Z_t = \frac{Cost}{Income}\times 100\%, \quad
TP = \frac{Income}{Cost}
\end{equation}

\subsection{Тооцооллын суурь өгөгдөл, таамаглал}
% Эдгээр нь таамаглал. Туршилтын байгууллагын бодит тоо (ажилтны тоо, цалин,
% хэмнэсэн цаг)-г авсны дараа шинэчилнэ.
\begin{itemize}
    \item Ашиглалтын хугацаа $n = 5$ жил, хорогдлын хувь $r = 12\%$.
    \item Ажилтны нэг цагийн дундаж өртөг 12\,000 төгрөг.
    \item Ажилтан бүр ажлын явц тодруулах, мэдээлэл хайх, тайлагнахад долоо хоногт 2 цаг хэмнэнэ: $25 \times 2 \times 48 = 2\,400$ цаг/жил.
    \item Менежер бүр ажил хянах, мэдээлэл нэгтгэхэд долоо хоногт 3 цаг хэмнэнэ: $5 \times 3 \times 48 = 720$ цаг/жил.
    \item Сарын тайлан бэлтгэх 24 цаг автоматжина: $24 \times 12 = 288$ цаг/жил.
    \item Нийт хэмнэлт 3\,408 цаг/жил; жилийн өгөөж $Income = 3\,408 \times 12\,000 = 40\,896\,000$ төгрөг.
    \item Жилийн зардал $Cost = 6\,900\,000$ төгрөг (засвар үйлчилгээ 6\,000\,000, сервер 600\,000, AI үйлчилгээ 300\,000).
    \item Жилийн цэвэр мөнгөн урсгал $Profit = Income - Cost = 33\,996\,000$ төгрөг.
\end{itemize}

\subsection{Анхны хөрөнгө оруулалтын зардал}
\begin{table}[H]
\centering
\caption{Анхны хөрөнгө оруулалтын бүтэц}
\label{table:inv}
\begin{tabular}{|p{8cm}|r|}
\hline
\textbf{Зардлын төрөл} & \textbf{Дүн (төгрөг)} \\ \hline
Судалгаа, системийн загварчлал & 3\,000\,000 \\ \hline
Backend, вэб, гар утасны хөгжүүлэлт (6 сар) & 18\,000\,000 \\ \hline
Тестлэлт, баримтжуулалт & 2\,000\,000 \\ \hline
Сервер, дэд бүтэц, апп дэлгүүрийн бүртгэл & 1\,000\,000 \\ \hline
\textbf{Нийт (Inv)} & \textbf{24\,000\,000} \\ \hline
\end{tabular}
\end{table}

\subsection{NPV тооцоолол}
\begin{table}[H]
\centering
\caption{NPV тооцооллын хүснэгт ($r = 12\%$, $FV_i = 33\,996\,000$ төгрөг)}
\label{table:npv}
\begin{tabular}{|c|r|c|r|}
\hline
\textbf{Жил} & \textbf{$FV_i$ (төгрөг)} & \textbf{$(1+r)^i$} & \textbf{$PV_i$ (төгрөг)} \\ \hline
1 & 33\,996\,000 & 1.12000 & 30\,353\,571 \\ \hline
2 & 33\,996\,000 & 1.25440 & 27\,101\,403 \\ \hline
3 & 33\,996\,000 & 1.40493 & 24\,197\,681 \\ \hline
4 & 33\,996\,000 & 1.57352 & 21\,605\,073 \\ \hline
5 & 33\,996\,000 & 1.76234 & 19\,290\,243 \\ \hline
\multicolumn{3}{|r|}{$\sum PV_i$} & 122\,547\,972 \\ \hline
\end{tabular}
\end{table}

\begin{equation}
NPV = 122\,547\,972 - 24\,000\,000 = 98\,547\,972 \text{ төгрөг}
\end{equation}
$NPV > 0$ тул төсөл эдийн засгийн хувьд үр ашигтай.

\subsection{Хөрөнгө оруулалтыг нөхөх хугацаа}
Эхний жилийн $PV_1 = 30\,353\,571$ төгрөг нь $Inv = 24\,000\,000$ төгрөгөөс их тул:
\begin{equation}
T_{inv} = 0 + \frac{24\,000\,000}{30\,353\,571} \approx 0.79 \text{ жил} \approx 9.5 \text{ сар}
\end{equation}

\subsection{Бусад үзүүлэлт}
\begin{equation}
PfN = \frac{122\,547\,972 / 5}{24\,000\,000} \approx 1.02, \qquad
P_a = \frac{33\,996\,000}{40\,896\,000}\times 100\% \approx 83.13\%
\end{equation}
\begin{equation}
P_c = \frac{33\,996\,000}{24\,000\,000} \approx 1.42, \qquad
Z_t = \frac{6\,900\,000}{40\,896\,000}\times 100\% \approx 16.87\%, \qquad
TP = \frac{40\,896\,000}{6\,900\,000} \approx 5.93
\end{equation}

\section{Эдийн засгийн үзүүлэлтүүдийн нэгдсэн дүн}
\begin{table}[H]
\centering
\caption{Эдийн засгийн үзүүлэлтүүдийн нэгдсэн дүн}
\label{table:econ}
\begin{tabular}{|p{8cm}|r|}
\hline
\textbf{Үзүүлэлт} & \textbf{Үр дүн} \\ \hline
Анхны хөрөнгө оруулалт (Inv) & 24\,000\,000 төгрөг \\ \hline
Жилийн өгөөж (хэмнэсэн цагийн үнэлгээ) & 40\,896\,000 төгрөг \\ \hline
Жилийн зардал & 6\,900\,000 төгрөг \\ \hline
Жилийн цэвэр мөнгөн урсгал & 33\,996\,000 төгрөг \\ \hline
NPV (5 жил, $r = 12\%$) & 98\,547\,972 төгрөг \\ \hline
Нөхөх хугацаа & $\approx$ 0.79 жил (9.5 сар) \\ \hline
Үр ашгийн дундаж норм (PfN) & 1.02 \\ \hline
Ашигт ажиллагаа ($P_a$) & 83.13\% \\ \hline
Зардлын түвшин ($Z_t$) & 16.87\% \\ \hline
Нийт бүтээмж (TP) & 5.93 \\ \hline
\end{tabular}
\end{table}

\section*{Бүлгийн дүгнэлт}
Тооцооллоор системийн NPV эерэг, хөрөнгө оруулалт нэг жилд хүрэхгүй хугацаанд нөхөгдөж, нэг төгрөгийн зардлаар 5.93 төгрөгийн өгөөж бий болох боломжтой байна. Тооцоо нь хэмнэгдэх ажлын цагийн таамаглалд суурилсан тул туршилтын ажиллагааны үед бодит хэмнэлтийг хэмжиж шинэчилнэ.
